\par\Needspace{4\baselineskip}
\begin{center}
\hypertarget{native-043-article027}{}\pdfbookmark[2]{27.}{article-027}27.
\end{center}
If equation IV is differentiated with $u$, $v$, and $\psi$ all treated as variables, we obtain
\[
\begin{aligned}
\frac{\mathrm{d}u}{u}&=\frac{\sin\psi\,\mathrm{d}v+\sin v\,\mathrm{d}\psi}{2\cos\frac{1}{2}(v-\psi)\cos\frac{1}{2}(v+\psi)}\\
&=\frac{r\operatorname{tang}\psi}{p}\,\mathrm{d}v+\frac{r\sin v}{p\cos\psi}\,\mathrm{d}\psi.
\end{aligned}
\]

Differentiating equation XI in the same way gives the following relation between the differential variations of $u$, $\psi$, and $N$
\SourceLeaf{044}
\[
\frac{\mathrm{d}N}{\lambda}=\left\{\frac{1}{2}e\left(1+\frac{1}{uu}\right)-\frac{1}{u}\right\}\mathrm{d}u+\frac{(uu-1)\sin\psi}{2u\cos\psi^2}\,\mathrm{d}\psi
\]
or
\[
\frac{\mathrm{d}N}{\lambda}=\frac{r}{bu}\,\mathrm{d}u+\frac{r\sin v}{b\cos\psi}\,\mathrm{d}\psi.
\]
Eliminating $\mathrm{d}u$ from this with the aid of the preceding equation, we obtain
\[
\frac{\mathrm{d}N}{\lambda}=\frac{rr}{bb\operatorname{tang}\psi}\,\mathrm{d}v+\left(1+\frac{r}{p}\right)\frac{r\sin v}{b\cos\psi}\,\mathrm{d}\psi
\]
or
\[
\begin{aligned}
\mathrm{d}v&=\frac{bb\operatorname{tang}\psi}{\lambda rr}\,\mathrm{d}N-\left(\frac{b}{r}+\frac{b}{p}\right)\frac{\sin v\operatorname{tang}\psi}{\cos\psi}\raisebox{-.6ex}{.}\,\mathrm{d}\psi\\
&=\frac{bb\operatorname{tang}\psi}{\lambda rr}\,\mathrm{d}N-\left(1+\frac{p}{r}\right)\frac{\sin v}{\sin\psi}\,\mathrm{d}\psi.
\end{aligned}
\]

\par\Needspace{4\baselineskip}
\begin{center}
\hypertarget{native-044-article028}{}\pdfbookmark[2]{28.}{article-028}28.
\end{center}
Differentiating equation X with $r$, $b$, $e$, and $u$ all treated as variables, substituting $\mathrm{d}e=\frac{\sin\psi}{\cos\psi^2}\,\mathrm{d}\psi$, and eliminating $\mathrm{d}u$ by means of the equation relating $\mathrm{d}N$, $\mathrm{d}u$, and $\mathrm{d}\psi$ given in the preceding article, we obtain
\[
\begin{aligned}
\mathrm{d}r&=\frac{r}{b}\,\mathrm{d}b+\frac{bbe(uu-1)}{2\lambda ur}\,\mathrm{d}N\\
&\quad+\frac{b}{2\cos\psi^2}\left\{\left(u+\frac{1}{u}\right)\sin\psi-\left(u-\frac{1}{u}\right)\sin v\right\}\mathrm{d}\psi.
\end{aligned}
\]
The coefficient of $\mathrm{d}N$ becomes $\frac{b\sin v}{\lambda\sin\psi}$ by equation VIII. In the coefficient of $\mathrm{d}\psi$, substituting $u(\sin\psi-\sin v)=\sin(\psi-v)$ and $\frac{1}{u}(\sin\psi+\sin v)=\sin(\psi+v)$ by equation IV gives $\frac{b\sin\psi\cos v}{\cos\psi^2}=\frac{p\cos v}{\sin\psi}$, so that we have
\[
\mathrm{d}r=\frac{r}{b}\,\mathrm{d}b+\frac{b\sin v}{\lambda\sin\psi}\,\mathrm{d}N+\frac{p\cos v}{\sin\psi}\,\mathrm{d}\psi.
\]
Further, when $N$ is regarded as a function of $b$ and $t$, we have $\mathrm{d}N=\frac{N}{t}\,\mathrm{d}t-\frac{3}{2}\cdot\frac{N}{b}\,\mathrm{d}b$. Substituting this value expresses $\mathrm{d}r$, and likewise $\mathrm{d}v$ in the preceding article, in terms of $\mathrm{d}t$, $\mathrm{d}b$, and $\mathrm{d}\psi$. We must also repeat here the earlier remark. To express the variations of the angles $v$ and $\psi$ in arcseconds in place of parts of the radius, we must either divide every term containing $\mathrm{d}v$ or $\mathrm{d}\psi$ by $206264{,}8$, or multiply all the remaining terms by this number.
